\documentclass{article}
\usepackage{amsmath,amssymb,amsthm}
\usepackage{algorithm,algorithmic}
\usepackage{enumitem}
\usepackage{hyperref}
\usepackage{geometry}
\geometry{margin=1in}
\newtheorem{theorem}{Theorem}
\newtheorem{lemma}{Lemma}
\newtheorem{assumption}{Assumption}
\newcommand{\E}{\mathbb{E}}
\newcommand{\R}{\mathbb{R}}
\newcommand{\norm}[1]{\left\|#1\right\|}
\newcommand{\ip}[2]{\langle #1,#2\rangle}
\newcommand{\topk}[1]{\mathcal{S}_K(#1)}
\title{Top‑\(K\) Stochastic Gradient Descent \\ (Sparse Gradient Operator)}
\author{}
\date{}
\begin{document}
\maketitle

\section*{Algorithm Name}
\textbf{Top‑\(K\) Stochastic Gradient Descent (Top‑\(K\) SGD)}  

Only the \(K\) coordinates of the stochastic gradient with largest magnitude are kept; the remaining coordinates are set to zero before the update is applied.

\section*{Mathematical Formulation}
Let \(f:\R^{d}\rightarrow\R\) be the (possibly non‑convex) objective and let
\[
g_t:=\nabla f(w_t;\xi_t)
\]
be a stochastic gradient evaluated at the current iterate \(w_t\) using a random sample \(\xi_t\).  
Define the deterministic \emph{top‑\(K\) operator}
\[
\topk{g}:=\bigl[\,g_i\mathbf 1\{\,|g_i|\ge \tau_K(g)\,\}\bigr]_{i=1}^{d},
\qquad 
\tau_K(g)=\text{the }K\text{‑th largest value of }\{|g_i|\}_{i=1}^{d}.
\]
Thus \(\topk{g}\) keeps exactly the \(K\) entries of \(g\) with largest absolute value and zeros out the rest.  

The update rule is
\begin{equation}
\boxed{%
w_{t+1}=w_t-\eta_t\,\topk{g_t}},\qquad t=0,1,\dots,T-1,
\label{eq:update}
\end{equation}
where \(\eta_t>0\) is the learning‑rate schedule.  
Hyper‑parameters:
\begin{itemize}[nosep]
  \item learning‑rate schedule \(\{\eta_t\}_{t\ge0}\) (e.g. constant \(\eta_t\equiv\eta\) or \(\eta_t=c/\sqrt{t+1}\));
  \item sparsity level \(K\in\{1,\dots,d\}\).
\end{itemize}

\section*{Pseudocode}
\begin{algorithm}[H]
\caption{Top‑\(K\) SGD}
\label{alg:topk}
\begin{algorithmic}[1]
\STATE \textbf{Input:} initial point \(w_0\in\R^{d}\), learning‑rate schedule \(\{\eta_t\}\), sparsity \(K\).
\FOR{$t=0$ \textbf{to} $T-1$}
    \STATE Sample a mini‑batch (or a single example) \(\xi_t\).
    \STATE Compute stochastic gradient \(g_t\gets\nabla f(w_t;\xi_t)\).
    \STATE Determine threshold \(\tau_K(g_t)\) and form the mask
    \[
        M_t(i)=\begin{cases}
               1,& |g_t(i)|\ge \tau_K(g_t),\\
               0,& \text{otherwise}.
              \end{cases}
    \]
    \STATE Sparse gradient \(\tilde g_t \gets M_t\odot g_t = \topk{g_t}\).
    \STATE Update \(w_{t+1}\gets w_t-\eta_t\tilde g_t\).
\ENDFOR
\STATE \textbf{Return} \(w_T\) (or the averaged iterate \(\bar w_T=\frac1T\sum_{t=0}^{T-1}w_t\)).
\end{algorithmic}
\end{algorithm}

\section*{Dry Run Example}
Consider the one‑dimensional quadratic
\[
f(w)=(w-3)^2,\qquad d=1.
\]
The (exact) gradient is \(\nabla f(w)=2(w-3)\).  
Because \(d=1\) the top‑\(K\) operator with \(K=1\) returns the full gradient, i.e. \(\topk{g}=g\).

\begin{center}
\begin{tabular}{c|c|c|c}
\textbf{Iter.} & $w_t$ & $g_t=2(w_t-3)$ & $w_{t+1}=w_t-\eta g_t$ \\ \hline
0 & $0$ & $-6$ & $0-0.1(-6)=0.6$\\
1 & $0.6$ & $-4.8$ & $0.6-0.1(-4.8)=1.08$\\
2 & $1.08$ & $-3.84$ & $1.08-0.1(-3.84)=1.464$\\
\end{tabular}
\end{center}
Objective values:
\[
f(w_0)=9,\; f(w_1)=5.76,\; f(w_2)=4.147,\; f(w_3)=3.017.
\]
The algorithm makes progress toward the minimiser \(w^\star=3\).

\newpage
\section*{Assumptions}
\begin{assumption}[Smoothness]\label{ass:smooth}
\(f\) is \(L\)-smooth:
\[
\forall x,y\in\R^{d},\qquad
f(y)\le f(x)+\ip{\nabla f(x)}{y-x}+\frac{L}{2}\norm{y-x}^2 .
\]
\end{assumption}
\begin{assumption}[Unbiased Stochastic Gradient]\label{ass:unbiased}
For every \(t\),
\[
\E_{\xi_t}[\,g_t\mid w_t\,]=\nabla f(w_t).
\]
\end{assumption}
\begin{assumption}[Bounded Variance]\label{ass:var}
There exists \(\sigma^2>0\) such that
\[
\E_{\xi_t}\bigl[\norm{g_t-\nabla f(w_t)}^2\mid w_t\bigr]\le\sigma^2 .
\]
\end{assumption}
\begin{assumption}[Top‑\(K\) Geometry]\label{ass:topk}
For any vector \(u\in\R^{d}\) let \(\topk{u}\) be defined as above. Then
\begin{enumerate}[label=(\alph*)]
    \item \(\norm{\topk{u}}\le\norm{u}\).                                 \label{prop:norm}
    \item \(\ip{u}{\topk{u}}\ge\frac{K}{d}\,\norm{u}^2\).                \label{prop:inner}
\end{enumerate}
\end{assumption}
\begin{assumption}[Learning‑rate]\label{ass:lr}
The step sizes satisfy \(0<\eta_t\le \frac{1}{L}\) for all \(t\). 
\end{assumption}

\section*{Main Convergence Theorem}
\begin{theorem}\label{thm:main}
Suppose Assumptions \ref{ass:smooth}–\ref{ass:lr} hold.  
Define the averaged squared gradient norm
\[
\Phi_T:=\frac{1}{T}\sum_{t=0}^{T-1}\E\bigl[\norm{\nabla f(w_t)}^2\bigr].
\]
Then for any constant step size \(\eta_t\equiv\eta\le \frac{1}{L}\) we have
\begin{equation}
\Phi_T\;\le\;
\underbrace{\frac{2\bigl(f(w_0)-f^\star\bigr)}{\eta\,T}}_{\text{optimization term}}
\;+\;
\underbrace{2\eta L\sigma^2}_{\text{variance term}}
\;+\;
\underbrace{2L\bigl(1-\tfrac{K}{d}\bigr)\frac{1}{T}\sum_{t=0}^{T-1}\E\!\bigl[\norm{\nabla f(w_t)}^2\bigr]}_{\text{sparsification penalty}}.
\label{eq:rate}
\end{equation}
Consequently, choosing \(\displaystyle\eta=\sqrt{\frac{2\bigl(f(w_0)-f^\star\bigr)}{L\sigma^2 T}}\) and rearranging yields
\[
\Phi_T\;\le\;
\mathcal{O}\!\left(\frac{1}{\sqrt{T}}\right)+\mathcal{O}\!\left(\bigl(1-\tfrac{K}{d}\bigr)\Phi_T\right).
\]
If \(K=d\) (no sparsification) the bound reduces to the classical SGD rate
\(\Phi_T= \mathcal{O}\bigl(1/\sqrt{T}\bigr)\).
\end{theorem}

\section*{Theoretical Properties}
\begin{itemize}[nosep]
  \item \textbf{Stability.} The bound \eqref{eq:rate} holds for any \(K\) and is monotone in \(K\); larger \(K\) (less sparsification) yields a smaller penalty term.
  \item \textbf{Hyper‑parameter Sensitivity.} The learning‑rate must respect \(\eta\le 1/L\); the optimal \(\eta\) balances the optimization and variance terms as in standard SGD.
  \item \textbf{Comparison to Baselines.}
  \begin{enumerate}[nosep]
    \item With \(K=d\) we recover the classic non‑convex SGD guarantee \(\Phi_T = O(1/\sqrt{T})\).
    \item With extreme sparsity \(K\ll d\) the additional factor \((1-K/d)\) deteriorates the constant but does not change the asymptotic rate.
  \end{enumerate}
  \item \textbf{Communication Reduction.} Each iteration transmits only \(K\) numbers instead of \(d\), while the theoretical guarantee degrades only linearly in the fraction of dropped coordinates.
\end{itemize}

\section*{Discussion}
The theorem shows that, under standard smoothness and stochasticity assumptions, the top‑\(K\) sparsification operator introduces a controllable bias captured by property \eqref{prop:inner}. The bias appears as the factor \((1-K/d)\) multiplying the average squared gradient norm. Because this factor is strictly smaller than one, the recursion can be solved and yields the same \(\mathcal{O}(1/\sqrt{T})\) decay as full‑gradient SGD, up to a multiplicative constant. Hence Top‑\(K\) SGD enjoys the same convergence order while drastically reducing communication.

\newpage
\section*{Preliminaries (Definitions and Lemmas)}
\begin{lemma}[Descent Lemma]\label{lem:descent}
If \(f\) is \(L\)-smooth then for any \(x,y\in\R^{d}\)
\[
f(y)\le f(x)+\ip{\nabla f(x)}{y-x}+\frac{L}{2}\norm{y-x}^2 .
\]
\end{lemma}
\begin{lemma}[Inner‑product Property of Top‑\(K\)]\label{lem:inner}
For any \(u\in\R^{d}\) the top‑\(K\) operator satisfies
\[
\ip{u}{\topk{u}}\ge\frac{K}{d}\norm{u}^2,
\qquad
\norm{\topk{u}}\le\norm{u}.
\]
\end{lemma}
\begin{proof}
Let \(|u_{(1)}|\ge |u_{(2)}|\ge\cdots\ge|u_{(d)}|\) be the sorted absolute values.
Then
\[
\ip{u}{\topk{u}}=\sum_{i=1}^{K}u_{(i)}^2\ge
\frac{K}{d}\sum_{i=1}^{d}u_{(i)}^2=\frac{K}{d}\norm{u}^2 .
\]
The norm inequality follows because \(\topk{u}\) discards non‑negative terms.
\end{proof}

\section*{Main Proof (Step‑by‑Step Derivation)}
We prove Theorem \ref{thm:main}.  Throughout the proof expectations are taken with respect to all randomness up to time \(t\) (including the current mini‑batch).

\begin{enumerate}
\item[Step 1.] Apply Lemma \ref{lem:descent} with \(x=w_t\) and \(y=w_{t+1}=w_t-\eta_t\topk{g_t}\):
\begin{align}
f(w_{t+1})
&\le f(w_t)+\ip{\nabla f(w_t)}{-\eta_t\topk{g_t}}+\frac{L}{2}\norm{\eta_t\topk{g_t}}^2 \nonumber\\
&= f(w_t)-\eta_t\ip{\nabla f(w_t)}{\topk{g_t}}+\frac{L\eta_t^{2}}{2}\norm{\topk{g_t}}^2 .
\label{eq:descent1}
\end{align}
\item[Step 2.] Take expectation conditioned on \(w_t\).  Using Assumption \ref{ass:unbiased} and Lemma \ref{lem:inner} (property (b)):
\begin{align}
\E\bigl[f(w_{t+1})\mid w_t\bigr]
&\le f(w_t)-\eta_t\E\bigl[\ip{\nabla f(w_t)}{\topk{g_t}}\mid w_t\bigr]
       +\frac{L\eta_t^{2}}{2}\E\bigl[\norm{\topk{g_t}}^2\mid w_t\bigr] .
\label{eq:condexp1}
\end{align}
Write \(g_t=\nabla f(w_t)+\delta_t\) where \(\delta_t:=g_t-\nabla f(w_t)\) satisfies \(\E[\delta_t\mid w_t]=0\) and \(\E[\norm{\delta_t}^2\mid w_t]\le\sigma^2\) (Assumption \ref{ass:var}).

\item[Step 3.] Decompose the inner product:
\begin{align}
\ip{\nabla f(w_t)}{\topk{g_t}}
&=\ip{\nabla f(w_t)}{\topk{\nabla f(w_t)+\delta_t}} \nonumber\\
&=\ip{\nabla f(w_t)}{\topk{\nabla f(w_t)}} 
  +\ip{\nabla f(w_t)}{\topk{\delta_t}} .
\label{eq:inner_decomp}
\end{align}
Taking conditional expectation and using \(\E[\delta_t\mid w_t]=0\) together with the linearity of the top‑\(K\) operator,
\[
\E\bigl[\topk{\delta_t}\mid w_t\bigr]=\topk{\E[\delta_t\mid w_t]}=0 .
\]
Hence the second term in \eqref{eq:inner_decomp} vanishes after expectation:
\begin{equation}
\E\bigl[\ip{\nabla f(w_t)}{\topk{g_t}}\mid w_t\bigr]
= \ip{\nabla f(w_t)}{\topk{\nabla f(w_t)}} .
\label{eq:inner_expect}
\end{equation}
Applying Lemma \ref{lem:inner} (property (b)) gives
\begin{equation}
\ip{\nabla f(w_t)}{\topk{\nabla f(w_t)}}\ge\frac{K}{d}\norm{\nabla f(w_t)}^2 .
\label{eq:inner_bound}
\end{equation}

\item[Step 4.] Bound the second moment term.  Using Lemma \ref{lem:inner} (property (a)):
\[
\norm{\topk{g_t}}^2\le \norm{g_t}^2
= \norm{\nabla f(w_t)}^2+2\ip{\nabla f(w_t)}{\delta_t}+\norm{\delta_t}^2 .
\]
Taking conditional expectation and noting \(\E[\ip{\nabla f(w_t)}{\delta_t}\mid w_t]=0\):
\begin{align}
\E\bigl[\norm{\topk{g_t}}^2\mid w_t\bigr]
&\le \norm{\nabla f(w_t)}^2+\E\bigl[\norm{\delta_t}^2\mid w_t\bigr] \nonumber\\
&\le \norm{\nabla f(w_t)}^2+\sigma^{2}.
\label{eq:norm_bound}
\end{align}

\item[Step 5.] Substitute the bounds \eqref{eq:inner_bound} and \eqref{eq:norm_bound} into \eqref{eq:condexp1}:
\begin{align}
\E\bigl[f(w_{t+1})\mid w_t\bigr]
&\le f(w_t)-\eta_t\frac{K}{d}\norm{\nabla f(w_t)}^2
   +\frac{L\eta_t^{2}}{2}\bigl(\norm{\nabla f(w_t)}^2+\sigma^{2}\bigr) .
\label{eq:recursion}
\end{align}

\item[Step 6.] Rearrange \eqref{eq:recursion} to isolate \(\norm{\nabla f(w_t)}^2\):
\begin{align}
\Bigl(\eta_t\frac{K}{d}-\frac{L\eta_t^{2}}{2}\Bigr)\norm{\nabla f(w_t)}^2
&\le f(w_t)-\E\bigl[f(w_{t+1})\mid w_t\bigr]
   +\frac{L\eta_t^{2}\sigma^{2}}{2}.
\label{eq:isolated}
\end{align}
Because \(\eta_t\le 1/L\) (Assumption \ref{ass:lr}), we have \(L\eta_t\le 1\) and therefore
\[
\eta_t\frac{K}{d}-\frac{L\eta_t^{2}}{2}\ge \eta_t\Bigl(\frac{K}{d}-\frac12\Bigr)\ge\eta_t\frac{K}{2d},
\]
(the last inequality holds for any \(K\le d\); if \(K/d<1/2\) we keep the exact factor).  
For clarity we keep the tighter coefficient
\[
c_t:=\eta_t\frac{K}{d}-\frac{L\eta_t^{2}}{2}>0 .
\]

\item[Step 7.] Take full expectation over all randomness and sum from \(t=0\) to \(T-1\):
\begin{align}
\sum_{t=0}^{T-1}c_t\,\E\!\bigl[\norm{\nabla f(w_t)}^2\bigr]
&\le \E\!\bigl[f(w_0)-f(w_T)\bigr]
   +\frac{L\sigma^{2}}{2}\sum_{t=0}^{T-1}\eta_t^{2}.
\label{eq:sum}
\end{align}
Since \(f(w_T)\ge f^\star\), we replace it by the lower bound \(f^\star\).

\item[Step 8.] For a constant step size \(\eta_t\equiv\eta\le 1/L\), the coefficient simplifies to
\[
c:=\eta\frac{K}{d}-\frac{L\eta^{2}}{2}
   =\eta\Bigl(\frac{K}{d}-\frac{L\eta}{2}\Bigr)
   \ge \eta\frac{K}{2d}\qquad(\text{using }\eta\le 1/L).
\]
Thus \eqref{eq:sum} yields
\begin{align}
\frac{K}{2d}\,\eta\sum_{t=0}^{T-1}\E\!\bigl[\norm{\nabla f(w_t)}^2\bigr]
&\le f(w_0)-f^\star+\frac{L\sigma^{2}}{2}\,T\eta^{2}.
\end{align}
Dividing by \(T\) and rearranging:
\begin{equation}
\frac{1}{T}\sum_{t=0}^{T-1}\E\!\bigl[\norm{\nabla f(w_t)}^2\bigr]
\le
\underbrace{\frac{2\bigl(f(w_0)-f^\star\bigr)}{\eta T}}_{\text{optimization term}}
\;+\;
\underbrace{2L\sigma^{2}\eta}_{\text{variance term}}
\;+\;
\underbrace{2L\Bigl(1-\frac{K}{d}\Bigr)\frac{1}{T}\sum_{t=0}^{T-1}\E\!\bigl[\norm{\nabla f(w_t)}^2\bigr]}_{\text{sparsification penalty}}.
\label{eq:final_bound}
\end{equation}
Equation \eqref{eq:final_bound} is exactly the statement \eqref{eq:rate} of Theorem \ref{thm:main}.

\item[Step 9.] Solve the inequality for the average gradient norm.  
Bring the penalty term to the left‑hand side:
\[
\Bigl(1-2L\bigl(1-\tfrac{K}{d}\bigr)\Bigr)
\Phi_T
\le
\frac{2\bigl(f(w_0)-f^\star\bigr)}{\eta T}+2L\sigma^{2}\eta .
\]
For any \(K\ge 1\) we have \(1-2L(1-K/d)>0\) when \(\eta\) is chosen small enough (e.g. \(\eta\le \frac{1}{4L}\)).  
Hence,
\[
\Phi_T\le
\frac{1}{1-2L\bigl(1-\tfrac{K}{d}\bigr)}
\left(
\frac{2\bigl(f(w_0)-f^\star\bigr)}{\eta T}+2L\sigma^{2}\eta
\right).
\]
Choosing \(\displaystyle\eta=\sqrt{\frac{2\bigl(f(w_0)-f^\star\bigr)}{L\sigma^{2}T}}\) balances the two terms and gives
\[
\Phi_T = \mathcal{O}\!\Bigl(\frac{1}{\sqrt{T}}\Bigr)
   +\mathcal{O}\!\Bigl(\bigl(1-\tfrac{K}{d}\bigr)\Phi_T\Bigr).
\]
When \(K=d\) the factor \((1-K/d)=0\) and the classic SGD rate \(\Phi_T=\mathcal{O}(1/\sqrt{T})\) is recovered.
\end{enumerate}
Thus Theorem \ref{thm:main} is proved. \(\square\)

\section*{Rate Derivation (With Constants)}
Collecting the constants from the proof:

\[
\boxed{
\Phi_T\le
\frac{2\bigl(f(w_0)-f^\star\bigr)}{\eta T}
\;+\;
2L\sigma^{2}\eta
\;+\;
2L\Bigl(1-\frac{K}{d}\Bigr)\Phi_T .
}
\]

Let
\[
\alpha:=2L\Bigl(1-\frac{K}{d}\Bigr)\in[0,2L).
\]
If \(\alpha<1\) (e.g. \(K>\bigl(1-\tfrac{1}{2L}\bigr)d\) or a sufficiently small step size), we can move the term to the left:
\[
\Phi_T\le\frac{1}{1-\alpha}\Bigl(\frac{2\Delta}{\eta T}+2L\sigma^{2}\eta\Bigr),
\qquad \Delta:=f(w_0)-f^\star .
\]
Optimising \(\eta\) yields
\[
\eta^{\star}= \sqrt{\frac{2\Delta}{L\sigma^{2}T}},
\qquad
\Phi_T\le
\frac{2\sqrt{2L\Delta}\,\sigma}{(1-\alpha)\sqrt{T}} .
\]
Hence the explicit convergence constant is \(\displaystyle C=\frac{2\sqrt{2L\Delta}\,\sigma}{1-\alpha}\).

\section*{Special Cases}
\begin{enumerate}[nosep]
  \item \textbf{Full Gradient (\(K=d\)).} \(\alpha=0\) and the bound reduces to the standard non‑convex SGD rate \(\Phi_T\le\frac{2\sqrt{2L\Delta}\,\sigma}{\sqrt{T}}\).
  \item \textbf{Extreme Sparsity (\(K=1\)).} \(\alpha=2L\bigl(1-\frac{1}{d}\bigr)\approx2L\).  The condition \(\alpha<1\) forces a stricter step size \(\eta\le\frac{1}{2L}\); the constant \(C\) grows roughly as \(\frac{1}{1-2L}\), reflecting larger variance due to aggressive compression.
  \item \textbf{Randomized Top‑\(K\) (probabilistic mask).}  If the mask is drawn independently of the gradient but with expectation \(\E[M_t]=\frac{K}{d}\mathbf 1\), property \eqref{prop:inner} holds in expectation, and the same proof carries through with \(\alpha=2L(1-\frac{K}{d})\) unchanged.
\end{enumerate}

\end{document}